\section{Ортогональные дополнения}
\begin{definition}
    Пусть $V$~--- векторное пространство, $U\subseteq V$. Имеется биекция:
    \[
        U \longleftrightarrow \operatorname{Ann} U \subseteq V^*.
    \]
    Если $V$~--- евклидово пространство, то также имеем:
    \[
        U \longleftrightarrow G_V^{-1}\operatorname{Ann} U \subseteq V.\\
    \]

    \[
    \operatorname{Ann} U = \{\varphi \in V^* \mid \varphi(u) = 0 \ \forall u \in U\}.
    \]

    Тогда
    \[
        G_V^{-1}\operatorname{Ann} U
        = \{v \in V \mid \langle u, v\rangle = 0 \ \forall u \in U\}
        =: U^\perp
    \]
    называется \emph{ортогональным дополнением} $U$.
\end{definition}

\begin{proplist}
    \item $\dim U^\perp = \dim \operatorname{Ann} U = \dim V - \dim U$.
    \item $(U^\perp)^\perp = U$.
    \item $U \cap U^\perp = \{0\}$
    \item
    \begin{enumerate}
        \item $(U + W)^\perp = U^\perp + W^\perp$
        \item $(U \cap W)^\perp = U^\perp \cap W^\perp$
    \end{enumerate}
    \begin{proof}
    \begin{proofsteps}
        \item
        $\dim U^\perp$ -- это размерность прообраза при изоморфизм аннулятора, а изморфизм не меняет размерность.
        \item
        Подсчитаем размерность:
		\[
			\dim(U^\perp)^\perp
			= \dim V - \dim U^\perp
			= \dim V - (\dim V - \dim U)
			= \dim U.
		\]
		Покажем включение $U \subset (U^\perp)^\perp$. Пусть $u \in U$ и $w \in U^\perp$. Тогда $\langle u, w\rangle = 0$,
  то есть $u \perp w$. Поскольку это верно для всех $w \in U^\perp$, получаем $u \in (U^\perp)^\perp$.

		Таким образом, $U \subset (U^\perp)^\perp$ и $\dim U = \dim(U^\perp)^\perp$, откуда $U = (U^\perp)^\perp$.
        \item
        Пусть $u \in U \cap U^\perp$. Тогда $u \perp U$, в частности $u \perp u$, то есть $\langle u, u\rangle = 0$,
        откуда $u = 0$.
        \item Упражнение

    \end{proofsteps}
    \end{proof}

\end{proplist}


\begin{cor}
	$V = U \oplus U^\perp = U + U^\perp $
	\begin{proof}
		По утверждению выше $U \cap U^\perp = \{0\}$, поэтому сумма прямая. По следствию из размерностей $\dim U + \dim U^\perp = \dim V$,
  откуда $U + U^\perp = V$.
	\end{proof}
\end{cor}

\begin{definition}
	Отображение $\pi_U : V \to U$ называется \emph{ортогональной проекцией} на $U$. Из разложения $V = U \oplus U^\perp$ следует,
 что $\forall v \in V$ единственным образом представляется как
	\[
		v = \pi_U v + (v - \pi_U v),
	\]
	где $\pi_U v \in U$ и $v - \pi_U v \in U^\perp$.
\end{definition}


